% ============================================================================
% ABSTRACT
% Fichier : preambule/abstract.tex
% ============================================================================

\chapter*{ABSTRACT}
\markboth{Abstract}{}
\addcontentsline{toc}{chapter}{Abstract}
\thispagestyle{preliminaire}

\begin{sloppypar}

Classical cryptography evaluates the security of a proof by the yardstick
of its robustness against computationally bounded adversaries. Conversely,
the institutional paradigm requires that the proof remain intelligible to
a non-specialist arbitrator in order to inform a decision. Operational
deployments reveal a recurrent pathology. A proof that is irreproachable
from the mathematical point of view may prove unusable to the agent tasked
with deciding. This structural tension between the two paradigms
constitutes the paradox of invisible authenticity, according to which the
more effectively a proof protects, the more opaque it becomes to a lay
third party.

The thesis introduces a novel informational metric, opposability Opp,
defined as the proportion of institutional uncertainty that a proof
succeeds in dissipating. It proposes the CLO framework, which aggregates
confidentiality (Conf), reliability (Rel) and opposability (Opp) and
articulates it with explicit cognitive parameters, in particular the
institutional extraction capacity $\Delta$ and the granularity of verdicts
$\alpha$. It establishes a conditional impossibility trilemma. Under four
hypotheses bounding the capacities of the institution, the product
Conf$\cdot$Rel$\cdot$Opp is bounded above by $\Delta/(\gamma\lambda)$,
where $\gamma = \alpha H(C) - 1$, up to a term of order $\lambda^{-2}$ and
a negligible remainder. A refined form identifies the regime in which institutional capacity
becomes the binding constraint. The limiting resource is then the quantity
of information that an institution can humanly extract from a proof, far
more than the size of the keys.

This result extends to the post-quantum setting. In the scenario of a
quantum support or verifier, the noise induces an interpretability loss $\eta_q$, which vanishes for
classical proofs and whose effect on the bound grows with the contextual
entropy. The thesis also proposes the SOTP family,
parameterised for different regimes of contextual richness, and
establishes its cryptographic properties. Its asymptotic optimality is
stated as a conjecture. The
algebraic core of the trilemma is mechanically certified in Lean4, which
makes explicit the hypotheses, constants and domain of validity and provides two formal witnesses guaranteeing the internal consistency and
the scope of the result.

Finally, the thesis specifies a refutable experimental programme
comprising six protocols intended to measure the institutional parameters,
with estimators and refutation thresholds, followed by an illustrative calibration of
these parameters on deployed systems. It checks the plausibility of the
orders of magnitude without constituting an empirical validation, since
the programme has not been executed under prospective conditions.

By raising opposability to the rank of a first-class cryptographic
property, this work constitutes institutional cryptography as an
autonomous theoretical domain. It provides instruments to guide the choice
of protocols in sensitive sectors such as electronic voting, regulatory
compliance and the post-quantum transition and prepares a future
standardisation integrating interpretability among the security criteria.

\vspace{0.5cm}
\noindent\textbf{Keywords:} Institutional cryptography, opposability,
CRO Trilemma, SOTP protocols, formal verification, falsifiable experimental
programme, Lean4.

\end{sloppypar}